\begin{afterword}
\summaryhead

The post continues the author's account of the years 1997 to 2002. By 2001 the young author had added Friendly AI to the old plan as a contingency, in case ``life turns out to be meaningless'', without admitting that the 1997 view had been a mistake. The contingency gave a line of retreat: the old view could be doubted without despair, but also given up a little at a time.

What was needed, the author says, was to stop and rethink everything built on the assumption that AI could not be dangerous. Instead the old arguments against regulating technology were rehearsed, and the old strategy stayed: build AI as fast as possible, start coding as soon as the money came, prefer openness to secrecy. Only the 1999 idea of an open-source AI project was dropped.

The lesson is that one large admission of error is more efficient than many small ones, since each small admission lets the old plan repair its justifications.

\responsehead

This is memoir, and as memoir it is candid. It describes the author's own slowness and does not excuse it. Where its account of the period can be checked against texts of the time, it holds up. \textsc{Creating Friendly AI} (2001) treats objective morality as ``one of the possibilities'' and says ``I simply do not believe the claim that relinquishment is possible,'' as the post says the young author did. The Institute's introduction page of the time, as quoted by Ben Garfinkel in 2022, expected its AI to reach ``transhumanity'' around 2008 or 2010, which fits ``gotta move now''. The aside on Szilard and Fermi is correct: they disagreed early on whether to keep fission research secret.

Some descriptions could not be checked, because the archived pages of 1996 to 2000 were out of reach: the strategy formed ``in the total absence'' of Friendly AI, the 1999 open-source plan, and the phrase ``the unlikely event that life turns out to be meaningless'', which I could not trace to a text of the period.

The lesson is not new to the reader of this book. It is the lesson of ``The Importance of Saying `Oops' '' (2007), whose own case was withheld; this sequence supplies that case. The general rule, that one large admission is ``far more efficient'' than many small ones, rests on that case, and the post adds no other evidence for it. It also sits beside ``Update Yourself Incrementally'' (2007), which advised shifting belief ``a little'' on weakly contrary evidence. The two fit if the difference is that the flaw here was already known, as the post says it was, but the post does not draw the line.

The remedy is named and not described. The post says a broken foundational assumption calls for stopping and rethinking everything, and that this is ``an art I need to write more about.'' The author did so in ``Crisis of Faith'' a little over two weeks later.

In short: a candid memoir whose account of 2001 matches the texts of the time where they could be checked, and whose lesson repeats ``The Importance of Saying `Oops' '' with the case that post withheld, without saying how it squares with updating incrementally.
\end{afterword}
